% TOP_PROBLEM: {"id": 8000003, "title": "Riemann-Hilbert Problem with non-analytic data", "category_id": 16, "category": {"id": 16, "name": "physics", "display_name": "Mathematical Physics", "description": "Problems at the intersection of mathematics and physics.", "slug": "physics", "order_index": 16, "created_at": "2026-07-31T15:26:25.671Z"}, "status": "partially_solved", "statusQualification": "This method-development problem leaves the minimal differentiability and error estimates to be determined. The source mentions Miller--McLaughlin and Varzugin work as partial context, not as a general solution."}
% FIRST_POSTED: 2026-10-11
% ENTRY_KIND: statement_only
% SUBMISSION: {"kind": "statement_only", "category": "top_problems", "problem_id": "8000003", "proposal_type": "established", "source": "Percy Deift, Some open problems in random matrix theory and the theory of integrable systems (2007), arXiv:0712.0849, Problem 3. https://arxiv.org/abs/0712.0849", "status_evidence": "Inherited source-catalogue status partially_solved preserved. Primary source accessed 11 October 2026; compared against the supplied catalogue excerpt. Historical wording and partial-result remarks were checked for scope; no comprehensive current-literature or proof-verification audit was performed.", "status_checked_on": "2026-10-11", "duplicate_check": "Full 15,773-record source registry checked for identical normalized statements; available site registry checked for identical normalized title/statement and canonical-ID collision. Canonical ID is absent from the site registry. No mathematical-equivalence literature audit performed.", "attribution": "Percy Deift. Statement preparation: OpenAI GPT-6 Codex assistant, generation 11 October 2026; source-matching and notation editing by the assistant, no independent human mathematical verification. Inherited rights/provenance note: Public arXiv source TeX; arXiv reuse license NEEDS_REVIEW"}
% SOURCE_STATEMENT: "In many situations one is concerned with the asymptotic behavior of Riemann-Hilbert problems with exponentially varying data of the form $e^{in\\phi(z)} r(z)$, $n \\to \\infty$. The Deift-Zhou nonlinear steepest descent method for such problems requires $\\phi(z)$ to be analytic. The analyticity is used in two ways: to control the equilibrium measure associated with the problem, and then to deform the contour for the RHP. By contrast, the method only requires minimal smoothness for $r(z)$ (see eg. Deift-Zhou in the context of Problem 12 below). It is of considerable theoretical and practical interest to extend the nonlinear steepest descent method to situations where $\\phi$ is no longer analytic, and has, for example, only a finite number of derivatives. For very interesting work on the analyticity problem, we refer the reader to a recent paper of Miller and McLaughlin. There is also interesting, older work due to Varzugin."
% Preparation and final self-audit: CONTRIBUTING.md prompts C and D; source/audit artifacts retained in the local physics_statements_2026_10_11 batch.
% Imported from: physics_statements.tex; UnsolvedMath ID 8000003
% !TeX program = lualatex
\documentclass[11pt]{article}
\usepackage[margin=25mm]{geometry}
\usepackage{fontspec}
\setmainfont{Latin Modern Roman}
\usepackage{amsmath,amssymb,mathtools}
\usepackage{unicode-math}
\setmathfont{Latin Modern Math}
\usepackage{xurl,hyperref}
\hypersetup{colorlinks=true,urlcolor=blue}
\setcounter{secnumdepth}{0}
\setlength{\parindent}{0pt}
\setlength{\parskip}{5pt}
\setlength{\emergencystretch}{3em}
\newcommand{\sref}[2]{\hyperlink{source:#1:#2}{[S#2]}}
\newcommand{\cataloguescope}[1]{\paragraph{Recorded target.}#1}
\newcommand{\R}{\mathbb R}
\newcommand{\C}{\mathbb C}
\newcommand{\N}{\mathbb N}
\newcommand{\Z}{\mathbb Z}
\newcommand{\Q}{\mathbb Q}
\newcommand{\F}{\mathbb F}
\newcommand{\eps}{\varepsilon}
\begin{document}
% UnsolvedMath ID 8000003; source catalogue code AMR-079-0003
\setcounter{section}{8000002}
\section{Riemann-Hilbert Problem with non-analytic data}\label{8000003}
{\small UnsolvedMath ID 8000003 · Mathematical Physics}
\subsection{Definitions and mathematical statement}

\paragraph{Shared notation.}$\mathbb N=\{1,2,\ldots\}$; the standard number fields have their usual meanings. All additional conventions are stated locally.


\paragraph{Definitions and notation.}A matrix Riemann--Hilbert problem asks for a matrix function analytic off an oriented contour, with prescribed jumps $M_+(z)=M_-(z)J_n(z)$ and normalization. The oscillatory data contain $e^{in\phi(z)}r(z)$, with $n\to\infty$. Nonlinear steepest descent deforms the contour and analyzes local and global model problems. The source asks to relax analyticity of the phase, both in equilibrium-measure analysis and contour deformation.

\paragraph{Catalogue wording retained for traceability.}
In many situations one is concerned with the asymptotic behavior of Riemann-Hilbert problems with exponentially varying data of the form $e^{in\phi(z)} r(z)$, $n \to \infty$. The Deift-Zhou nonlinear steepest descent method for such problems requires $\phi(z)$ to be analytic. The analyticity is used in two ways: to control the equilibrium measure associated with the problem, and then to deform the contour for the RHP. By contrast, the method only requires minimal smoothness for $r(z)$ (see eg. Deift-Zhou in the context of Problem 12 below). It is of considerable theoretical and practical interest to extend the nonlinear steepest descent method to situations where $\phi$ is no longer analytic, and has, for example, only a finite number of derivatives. For very interesting work on the analyticity problem, we refer the reader to a recent paper of Miller and McLaughlin. There is also interesting, older work due to Varzugin.
\paragraph{Scope and source qualification.}This method-development problem leaves the minimal differentiability and error estimates to be determined. The source mentions Miller--McLaughlin and Varzugin work as partial context, not as a general solution.
\subsection{Short English statement}
Extend nonlinear steepest descent for oscillatory Riemann--Hilbert problems to phases with only finitely many derivatives.
\subsection{Sources}
\begin{itemize}
\item[\hypertarget{source:8000003:1}{S1}] ULAM.AI and the UnsolvedMath Contributors, UnsolvedMath, record 8000003 (AMR-079-0003). Catalogue curation under CC BY 4.0; definitions and English exposition added locally.
\url{https://huggingface.co/datasets/ulamai/UnsolvedMath}.
\item[\hypertarget{source:8000003:2}{S2}] Percy Deift, Some open problems in random matrix theory and the theory of integrable systems (2007), arXiv:0712.0849, Problem 3.
\url{https://arxiv.org/abs/0712.0849}.
\end{itemize}
\label{end:8000003}
\end{document}
